% 第5章 Heisenberg 模型的基态
\chapter{Heisenberg 模型的基态}

3.2 节中，自旋 1/2 的反铁磁 Heisenberg 模型作为 Mott 绝缘体的有效哈密顿量出现。一般地说，Heisenberg 哈密顿量是量子磁学的基本模型，也是其他可以用量子自旋算符有效描述的现象的基本模型\footnote{例如 3.3 节所示的超导电性与电荷密度波。}。从其基态、激发与热力学相的研究中，可以学到极宽范围内的一批概念和技术。

大小为 $S$ 的 $\mathcal{N}$ 个自旋的格子，其哈密顿量为

\begin{align*}
\mathcal {H} & = \mathcal {H} ^ {zz} + \mathcal {H} ^ {xy} ,\\
\mathcal {H} ^ {zz} & = \frac {1}{2} \sum_ {ij} J _ {ij} S _ {i} ^ {z} S _ {j} ^ {z} ,\\
\mathcal {H} ^ {xy} & = \frac {1}{4} \sum_ {ij} J _ {ij} \left( S _ {i} ^ {+} S _ {j} ^ {-} + S _ {i} ^ {-} S _ {j} ^ {+} \right) ,\\
S _ {\mathrm{tot}} ^ {z} & = \sum_ {i} S _ {i} ^ {z} ,
\tag{5.1}
\end{align*}

其中 $J_{ij} = J_{ji}$ 具有格子平移对称性。$\mathcal{H}$ 与总自旋

\begin{equation*}
\mathbf {S} _ {\mathrm{tot}} = \sum_ {i} \mathbf {S} _ {i}
\tag{5.2}
\end{equation*}

的所有三个分量对易，因而是转动不变的。于是本征态可标记为

\begin{align*}
\Psi & = | S _ {\mathrm{tot}}, M, \ldots \rangle ,\\
M & = - S _ {\mathrm{tot}}, - S _ {\mathrm{tot}} + 1, \ldots , S _ {\mathrm{tot}} ,\\
S _ {\mathrm{tot}} & \leq \mathcal {N} S ,
\tag{5.3}
\end{align*}

$M$ 是总磁化 $S_{\mathrm{tot}}^{z}$ 的本征值，而

\begin{equation*}
\mathbf {S} _ {\mathrm{tot}} ^ {2} = S _ {\mathrm{tot}} \left( S _ {\mathrm{tot}} + 1 \right).
\tag{5.4}
\end{equation*}

若耦合 $J_{ij}$ 对 $i$、$j$ 的格子平移不变，则格子动量 $\mathbf{k}$ 也是好量子数。

本章集中于 Heisenberg 模型基态的两个方面。5.1 节证明关于总自旋与波函数符号的 Marshall 定理；5.2 节证明半奇整数自旋链在大 $\mathcal{N}$ 极限下的无能隙行为。

\section{反铁磁体}

一般而言，反铁磁基态远比铁磁基态复杂。下面的定理对二分反铁磁哈密顿量（见定义 3.1）的基态给出强约束。子晶格 A、B 上的交错磁化算符为

\begin{equation*}
S ^ {\mathrm{stagg}} = \sum_ {i \in A} S _ {i} ^ {z} - \sum_ {i \in B} S _ {i} ^ {z}.
\tag{5.5}
\end{equation*}

Ising 组态构成基底

\begin{equation*}
\Phi_ {\alpha} = \left| S, m _ {1} ^ {\alpha} \right\rangle_ {1} \left| S, m _ {2} ^ {\alpha} \right\rangle_ {2} \dots \left| S, m _ {\mathcal {N}} ^ {\alpha} \right\rangle_ {\mathcal {N}} ,
\tag{5.6}
\end{equation*}

其中 $|S, m_i \rangle_i$ 表示 $\mathbf{S}_i^2$、$S_i^z$ 的本征态，本征值分别为 $S(S + 1)$、$m_i$。使 $S^{\mathrm{stagg}}$ 极大的 Néel 态是 Ising 组态

\begin{equation*}
\Psi^ {\mathrm{N\acute{e}el}} = \prod_ {i} \left| S, \eta_ {i} S \right\rangle_ {i} \, , \quad \eta_ {i} = \left\{ \begin{array}{ll} 1 & i \in A \\ - 1 & i \in B \end{array} \right. .
\tag{5.7}
\end{equation*}

然而一般 $\Psi^{\mathrm{N\acute{e}el}}$ 并非 $\mathcal{H}$ 的本征态。“捣乱者”是 $H^{xy}$ 中的自旋翻转项 $S_{i}^{+}S_{j}^{-}$，它们把 $\Psi^{\mathrm{N\acute{e}el}}$ 与其他 Ising 组态联系起来。为了我们的目的，方便的做法是把子晶格 B 的自旋轴绕 $z$ 轴转动，即作幺正变换

\begin{align*}
i & \in B :\\
S _ {i} ^ {+} & \rightarrow - \tilde {S} _ {i} ^ {+} ,\\
S _ {i} ^ {-} & \rightarrow - \tilde {S} _ {i} ^ {-} ,\\
S _ {i} ^ {z} & \rightarrow + \tilde {S} _ {i} ^ {z} ,
\tag{5.8}
\end{align*}

它把 Ising 组态变为

\begin{equation*}
\left| S, m _ {i} \right\rangle_ {i} \to \left| S, \tilde {m} _ {i} \right\rangle_ {i} = \left\{ \begin{array}{ll} \left| S, m _ {i} \right\rangle_ {i} & i \in A \\ (- 1)^{(S + m _ {i})} \left| S, m _ {i} \right\rangle_ {i} & i \in B \end{array} \right. ,
\tag{5.9}
\end{equation*}

其中 $\tilde{m}_{i}$ 分别是子晶格 A 上 $S_{i}^{z}$ 与子晶格 B 上 $\tilde{S}_{i}^{z}$ 的本征值。把自己限制在某个 $M$ 子空间，把所有波函数展开为

\begin{equation*}
\Psi^ {M} = \sum_ {\alpha} f _ {\alpha} ^ {M} \tilde {\Phi} _ {\alpha} ^ {M},
\tag{5.10}
\end{equation*}

其中

\begin{equation*}
\tilde {\Phi} _ {\alpha} ^ {M} = \prod_{\sum_i \tilde m_i^\alpha = M} \left| S, \tilde {m} _ {i} ^ {\alpha} \right\rangle_ {i}.
\tag{5.11}
\end{equation*}

现在给出 Marshall 定理（经 Lieb 与 Mattis 推广，见文献注释）。

\medskip
\noindent\textbf{Marshall 定理 5.1}\quad 考虑哈密顿量 (5.1)，具有连接子晶格 $A$、$B$（二分）的反铁磁交换 $J_{ij} \geq 0$，且格子上任意两格点都可以经由中间格点的有限交换序列相连。则在任何允许的 $M$ 扇区中，最低能量态 $\Psi_0^M$ 可以取为在转动的 Ising 基 $\tilde{\Phi}_{\alpha}^{M}$ 中系数正定：

\begin{equation*}
\Psi_ {0} ^ {M} = \sum_ {\alpha} f _ {\alpha} ^ {M} \tilde {\Phi} _ {\alpha} ^ {M} \, , \quad f _ {\alpha} ^ {M} > 0, \; \forall \alpha .
\tag{5.12}
\end{equation*}
\medskip

因此，由 (5.8)，$\Psi_0^M$ 在未转动 Ising 组态 $|\Phi_{\alpha}\rangle$ 下的系数满足 Marshall 符号判据

\begin{align*}
\Psi_ {0} ^ {M} & = \sum_ {\alpha} ( - 1 ) ^ {\Gamma ( \alpha )} f _ {\alpha} ^ {M} \left| \Phi_ {\alpha} \right\rangle ,\\
\Gamma_ {\alpha} & = \sum_ {i \in B} \left( S + m _ {i} ^ {\alpha} \right) .
\tag{5.13}
\end{align*}

\medskip
\noindent\textbf{Marshall 定理 5.2}\quad 对等大小子晶格 $A$、$B$，绝对基态 $\Psi_0$ 是总自旋单态：

\begin{equation*}
\mathbf {S} _ {\mathrm{tot}} \left| \Psi_ {0} \right\rangle = 0.
\tag{5.14}
\end{equation*}
\medskip

首先必须强调：并非所有总自旋单态都满足 Marshall 符号判据 (5.12)；反过来，也不是所有满足 Marshall 符号的态都是总自旋单态。但我们这个 Heisenberg 反铁磁体的基态必须同时满足两个条件。

\textbf{证明：}用子晶格转动后的算符 (5.8)，哈密顿量 (5.1) 变为

\begin{align*}
\mathcal {H} & = \tilde {\mathcal {H}} ^ {zz} + \tilde {\mathcal {H}} ^ {xy} ,\\
\tilde {\mathcal {H}} ^ {zz} & = + \sum_ {i \in A,\, j \in B} | J _ {ij} | \, S _ {i} ^ {z} \tilde {S} _ {j} ^ {z} ,\\
\tilde {\mathcal {H}} ^ {xy} & = - \frac {1}{2} \sum_ {i \in A,\, j \in B} | J _ {ij} | \left( S _ {i} ^ {+} \tilde {S} _ {j} ^ {-} + S _ {i} ^ {-} \tilde {S} _ {j} ^ {+} \right) .
\tag{5.15}
\end{align*}

注意 $\tilde{H}^{zz}$ 在 Ising 组态中是对角的：

\begin{equation*}
\tilde {\mathcal {H}} ^ {zz} \left| \tilde {\Phi} _ {\alpha} \right\rangle = e _ {\alpha} \left| \tilde {\Phi} _ {\alpha} \right\rangle ,
\tag{5.16}
\end{equation*}

（上标 $M$ 暂时从 $\tilde{\Phi}^{M}$ 上略去，需要时再恢复。）关键之点是：在这个子晶格转动表象中，$\tilde{H}^{xy}$ 只有非正的矩阵元：

\begin{equation*}
\left\langle \tilde {\Phi} _ {\alpha} \left| \tilde {\mathcal {H}} ^ {xy} \right| \tilde {\Phi} _ {\beta} \right\rangle = - \left| K _ {\alpha \beta} \right| .
\tag{5.17}
\end{equation*}

系数 $f$ 的本征值方程为

\begin{equation*}
- \sum_ {\beta} | K _ {\alpha \beta} | \, f _ {\beta} + e _ {\alpha} f _ {\alpha} = E f _ {\alpha}.
\tag{5.18}
\end{equation*}

考虑试探函数

\begin{equation*}
\bar {\Psi} = \sum_ {\alpha} \left| f _ {\alpha} \right| \left| \tilde {\Phi} _ {\alpha} \right\rangle ,
\tag{5.19}
\end{equation*}

其能量为

\begin{align*}
\langle \bar {\Psi} | \tilde {\mathcal {H}} | \bar {\Psi} \rangle & = \sum_ {\alpha} e _ {\alpha} \left| f _ {\alpha} \right| ^ {2} - \sum_ {\alpha \beta} \left| K _ {\alpha \beta} \right| \left| f _ {\alpha} \right| \left| f _ {\beta} \right|\\
& \leq \sum_ {\alpha} e _ {\alpha} \left| f _ {\alpha} \right| ^ {2} - \sum_ {\alpha \beta} \left| K _ {\alpha \beta} \right| f _ {\alpha} f _ {\beta} = E.
\tag{5.20}
\end{align*}

因此它也必须是基态，并满足本征值方程

\begin{equation*}
- \sum_ {\beta} \left| K _ {\alpha \beta} \right| \left| f _ {\beta} \right| + e _ {\alpha} \left| f _ {\alpha} \right| = E \left| f _ {\alpha} \right| .
\tag{5.21}
\end{equation*}

由于 $\tilde{\Phi}_{\alpha}$ 并非 $\tilde{\mathcal{H}}$ 的本征态\footnote{除 $M = \mathcal{N}S$ 的情形——那是一个单态空间。}，$e_{\alpha}$ 大于其最低能量，即

\begin{equation*}
\forall \alpha \, , \quad e _ {\alpha} - E > 0.
\tag{5.22}
\end{equation*}

把 (5.18) 与 (5.21) 改写为

\begin{equation*}
\left( e _ {\alpha} - E \right) f _ {\alpha} = \sum_ {\beta} \left| K _ {\alpha \beta} \right| f _ {\beta},
\tag{5.23}
\end{equation*}

\begin{equation*}
\left( e _ {\alpha} - E \right) \left| f _ {\alpha} \right| = \sum_ {\beta} \left| K _ {\alpha \beta} \right| \left| f _ {\beta} \right| ,
\tag{5.24}
\end{equation*}

对两式两边取绝对值，得

\begin{equation*}
\left| \sum_ {\beta} \left| K _ {\alpha \beta} \right| f _ {\beta} \right| = \sum_ {\beta} \left| K _ {\alpha \beta} \right| \left| f _ {\beta} \right| ,
\tag{5.25}
\end{equation*}

这意味着

\begin{equation*}
f _ {\beta} \geq 0.
\tag{5.26}
\end{equation*}

于是试探态与基态是同一个态：$\bar{\Psi} = \pm\Psi$。还可以如下证明 $f_{\beta}$ 的更强条件。

\subsection*{引理 5.3}

\begin{equation*}
\forall \beta \, , \quad f _ {\beta} > 0.
\tag{5.27}
\end{equation*}

\textbf{证明：}容易验证，任何总磁化为 $M$ 的 Ising 组态，通过逐次应用成对自旋翻转算符 $K_{\alpha\beta}$，与同一扇区中的所有其他组态相连。于是若某个 $f_{\alpha}$ 为零，则由 (5.22) 与 (5.23)，同一 $M$ 扇区中所有 $f_{\beta}$ 也必须为零。既然至少要有一个系数非零，所有 $f_{\beta}$ 必须非零。这证明了引理 5.3，并完成定理 5.1 的证明。证毕。

\medskip
\noindent\textbf{推论 5.4}\quad 对任何固定的 $M$，$\Psi_{0}$ 非简并。
\medskip

这由 (5.27) 得出：在 $M$ 子空间中不可能存在与 $\Psi$ 正交而系数全正定的态。

为了证明定理 5.2，先要证明 $M$ 扇区基态具有最小可能的总自旋。

\subsection*{引理 5.5}

\begin{equation*}
\left( \mathbf {S} _ {\mathrm{tot}} \right) ^ {2} \left| \Psi^ {M} \right\rangle = M ( M + 1 ) \left| \Psi^ {M} \right\rangle .
\tag{5.28}
\end{equation*}

\textbf{证明：}考察两子晶格格点数相等的二分格子上的无穷程哈密顿量：

\begin{equation*}
\mathcal {H} ^ {\infty} = J \sum_ {i \in A,\, j \in B} \mathbf {S} _ {i} \cdot \mathbf {S} _ {j} = J \, \mathbf {S} _ {\mathrm{tot},A} \cdot \mathbf {S} _ {\mathrm{tot},B}.
\tag{5.29}
\end{equation*}

利用

\begin{equation*}
\mathbf {S} _ {\mathrm{tot}} = \mathbf {S} _ {\mathrm{tot},A} + \mathbf {S} _ {\mathrm{tot},B}
\tag{5.30}
\end{equation*}

这个模型作为两自旋问题可以平凡求解。总自旋算符的可能取值为

\begin{equation*}
S _ {\mathrm{tot},A} \, , S _ {\mathrm{tot},B} = 0, 1, \dots , \mathcal {N} S / 2 ,
\tag{5.31}
\end{equation*}

对等大小子晶格这意味着

\begin{equation*}
0 \leq S _ {\mathrm{tot}} \leq \mathcal {N} S ,
\tag{5.32}
\end{equation*}

而 (5.29) 的本征值为

\begin{equation*}
E ^ {\infty} \left( S _ {\mathrm{tot}} \right) = \frac {J}{2} \left[ S _ {\mathrm{tot}} \left( S _ {\mathrm{tot}} + 1 \right) - S _ {\mathrm{tot},A} \left( S _ {\mathrm{tot},A} + 1 \right) - S _ {\mathrm{tot},B} \left( S _ {\mathrm{tot},B} + 1 \right) \right].
\tag{5.33}
\end{equation*}

由于 (5.33) 随 $S_{tot}$ 单调增加，且 $S_{tot} > M$，给定 $M$ 扇区中的基态有

\begin{equation*}
S _ {\mathrm{tot}} = M.
\tag{5.34}
\end{equation*}

$\mathcal{H}$ 与 $\mathcal{H}^{\infty}$ 都满足 Marshall 定理 (5.1) 的要求，故其基态都有 Marshall 符号 (5.12)。因此它们的交叠（正数之和）不会为零，二者必须有相同的总自旋量子数。于是我们证明了 $\mathcal{H}$ 在扇区 $M$ 中的最低能量态自旋为 $S_{tot} = M$。既然总自旋的允许值 $S_{tot}' > M$ 在磁化为 $M$ 的扇区中都有代表，$E(S_{tot})$ 满足

\begin{equation*}
\forall S _ {\mathrm{tot}} ' > S _ {\mathrm{tot}} \; \Rightarrow \; E ( S _ {\mathrm{tot}} ' ) > E ( S _ {\mathrm{tot}} ),
\tag{5.35}
\end{equation*}

即能量是 $S_{tot}$ 的单调增函数。现在考虑 $M = 0$ 扇区。不等式 (5.35) 证明基态必须取最小可能的 $S_{tot}$，由 (5.32) 这是 $S_{tot} = 0$，定理 5.2 得证。证毕。

如 4.2 节所论，二分格子上 Heisenberg 反铁磁体的这种非负性，同样为半满的 Hubbard 模型和一切填充的负 $U$ Hubbard 模型所共有。它也是所有格子上 Heisenberg 铁磁体的性质。

\medskip
\noindent\textbf{推论 5.6}\quad 对非正交换的铁磁模型

\begin{equation*}
\mathcal {H} = - \frac {1}{2} \sum_ {ij} | J _ {ij} | \, \mathbf {S} _ {i} \cdot \mathbf {S} _ {j} ,
\tag{5.36}
\end{equation*}

完全铁磁态 $\Psi^{\mathrm{FM}}$ 属于基态多重态，其中

\begin{equation*}
\Psi^ {\mathrm{FM}} = \prod_ {i = 1} ^ {\mathcal {N}} \left| S, S \right\rangle_ {i}.
\tag{5.37}
\end{equation*}
\medskip

（证明留作习题，紧随反铁磁情形的证明。）

作为简并多重态的成员，$\Psi^{\mathrm{FM}}$ 自发破缺哈密顿量的转动对称性。这种对称破缺很特殊：算符 $S_{tot}^{z}$ 与哈密顿量对易，$M$ 是“好量子数”（即标记 $\mathcal{H}$ 的本征态）。在这方面反铁磁体不同于铁磁体：交错磁化一般与哈密顿量不对易。然而，在严格热力学极限（$\mathcal{N} = \infty$）下，反铁磁体仍然可能发生自发对称破缺，这将在 6.1 节讨论。

\section{半奇整数自旋链}

半奇整数自旋的粒子（如电子）具有一个奇特的量子性质：绕任何自旋轴转动 $2\pi$，其波函数获得一个负号。当波函数在自旋空间中局域于某个特定方向附近时，这一效应几乎不可察觉。但在低维，零点涨落很大——第 11 章的半经典自旋波理论将展示这一点。于是，这些负号引起的干涉效应可能变得重要。例如我们将看到，一维 Heisenberg 反铁磁体对整数与半奇整数自旋具有定性不同的谱。Lieb、Schultz 与 Mattis 的如下定理（见文献注释）确立了这一差别。

\medskip
\noindent\textbf{Lieb--Schultz--Mattis 定理 5.7}\quad 考虑链上的 Heisenberg 反铁磁体（$J > 0$）

\begin{equation*}
\mathcal {H} = J \sum_ {j = 1} ^ {\mathcal {N}} \mathbf {S} _ {j} \cdot \mathbf {S} _ {j + 1} + J \, \mathbf {S} _ {\mathcal {N}} \cdot \mathbf {S} _ {1} ,
\tag{5.38}
\end{equation*}

$\mathcal{N}$ 为偶数。对半奇整数自旋，如 $S = 1/2, 3/2, \ldots$，存在一个激发态，其能量随 $\mathcal{N} \to \infty$ 趋于零。
\medskip

$\Psi_0$ 是 $\mathcal{H}$ 的基态。考虑扭转算符

\begin{equation*}
\mathcal {O} ^ {1} = \exp \left( i \frac {2 \pi}{\mathcal {N}} \sum_ {j = 1} ^ {\mathcal {N}} j \, S _ {j} ^ {z} \right),
\tag{5.39}
\end{equation*}

它在 $x$--$y$ 平面内逐格点增量地转动每个自旋，使得从第一个格点到最末格点，自旋坐标绕 $z$ 轴共转动 $2\pi$。“扭转态”构造为

\begin{equation*}
\Psi_ {1} \equiv \mathcal {O} ^ {1} \left| \Psi_ {0} \right\rangle .
\tag{5.40}
\end{equation*}

证明由两个命题组成：

\subsection*{命题 5.8}

\begin{equation*}
\langle \Psi_ {1} | \Psi_ {0} \rangle = 0
\tag{5.41}
\end{equation*}

以及

\subsection*{命题 5.9}

\begin{equation*}
\lim _ {\mathcal {N} \to \infty} \left[ \langle \Psi_ {1} | \mathcal {H} | \Psi_ {1} \rangle - E _ {0} \right] \to 0 ,
\tag{5.42}
\end{equation*}

其中 $E_0$ 是 $\Psi_0$ 的能量。

命题 5.8 的证明基于 $\Psi_{1}$ 与 $\Psi_{0}$ 的格子动量相差 $\pi$，因而正交。于是把 $\Psi_{1}$ 按精确本征态展开（其中不含 $\Psi_{0}$），再用命题 5.9 即可证明：至少存在一个激发态，其激发能在大幅子极限下趋于零。

格子平移算符 $U_{x}$ 把自旋平移一个格点常量：

\begin{align*}
U _ {x} \mathbf {S} _ {j} U _ {x} ^ {- 1} & = \mathbf {S} _ {j + 1} ,\\
U _ {x} \mathbf {S} _ {\mathcal {N}} U _ {x} ^ {- 1} & = \mathbf {S} _ {1} .
\tag{5.43}
\end{align*}

哈密顿量平移不变：

\begin{equation*}
[ \mathcal {H}, U _ {x} ] = 0 ,
\tag{5.44}
\end{equation*}

且由推论 5.4，$\Psi_0$ 非简并。于是 $\Psi_0$ 必是 $U_x$ 的本征函数：

\begin{equation*}
U _ {x} \Psi_ {0} = e ^ {i k _ {0}} \Psi_ {0}.
\tag{5.45}
\end{equation*}

$k_{0}$ 是基态格子动量。扭转态与基态的交叠为

\begin{align*}
\langle \Psi_ {0} | \Psi_ {1} \rangle & = \langle \Psi_ {0} | \mathcal {O} ^ {1} \Psi_ {0} \rangle\\
& = \langle \Psi_ {0} | U _ {x} \mathcal {O} ^ {1} U _ {x} ^ {- 1} | \Psi_ {0} \rangle .
\tag{5.46}
\end{align*}

然而

\begin{equation*}
U _ {x} \mathcal {O} ^ {1} U _ {x} ^ {- 1} = \mathcal {O} ^ {1} \exp ( i 2 \pi S _ {1} ^ {z} ) \exp \left( - i \frac {2 \pi}{\mathcal {N}} S _ {\mathrm{tot}} ^ {z} \right).
\tag{5.47}
\end{equation*}

由 Marshall 定理 5.2，$\Psi_0$ 是总自旋单态，故

\begin{equation*}
\exp \left( - i \frac {2 \pi}{\mathcal {N}} S _ {\mathrm{tot}} ^ {z} \right) \left| \Psi_ {0} \right\rangle = \left| \Psi_ {0} \right\rangle .
\tag{5.48}
\end{equation*}

再用恒等式

\begin{equation*}
\exp \left( i 2 \pi S _ {1} ^ {z} \right) = \left\{ \begin{array}{ll} 1 & S = 0, 1, 2 \ldots \\ - 1 & S = \frac {1}{2}, \frac {3}{2}, \frac {5}{2}, \ldots \end{array} \right. ,
\tag{5.49}
\end{equation*}

于是 (5.47) 给出

\begin{equation*}
U _ {x} \mathcal {O} ^ {1} U _ {x} ^ {- 1} = \left\{ \begin{array}{ll} \mathcal {O} ^ {1} & S = 0, 1, 2 \ldots \\ - \mathcal {O} ^ {1} & S = \frac {1}{2}, \frac {3}{2}, \frac {5}{2}, \ldots \end{array} \right. .
\tag{5.50}
\end{equation*}

因此对半奇整数 $S$ 有

\begin{equation*}
\langle \Psi_ {0} | \Psi_ {1} \rangle = - \langle \Psi_ {0} | \Psi_ {1} \rangle = 0 ,
\tag{5.51}
\end{equation*}

命题 5.8 得证。证毕。

$\mathcal{O}^1$ 幺正，故 $\Psi_{1}$ 已归一化。扭转态的能量为

\begin{align*}
\langle \Psi_ {1} | \mathcal {H} | \Psi_ {1} \rangle & = E _ {0} + \Big\langle \Psi_ {0} \Big| \left[ \cos ( 2 \pi / \mathcal {N} ) - 1 \right] \sum_ {j = 1} ^ {\mathcal {N}} \left( S _ {j} ^ {x} S _ {j + 1} ^ {x} + S _ {j} ^ {y} S _ {j + 1} ^ {y} \right)\\
& \quad + \sin ( 2 \pi / \mathcal {N} ) \sum_ {j = 1} ^ {\mathcal {N}} \left( S _ {j} ^ {x} S _ {j + 1} ^ {y} - S _ {j} ^ {y} S _ {j + 1} ^ {x} \right) \Big| \Psi_ {0} \Big\rangle .
\tag{5.52}
\end{align*}

由于 $\Psi_{0}$ 是总自旋单态、转动不变，任何在转动下为奇的算符其期望值为零。特别地，把 $S^{x} \rightarrow S^{y}$、$S^{y} \rightarrow S^{x}$ 的整体转动用于上式给出

\begin{equation*}
\left\langle \Psi_ {0} \left| \sum_ {j = 1} ^ {\mathcal {N}} \left( S _ {j} ^ {x} S _ {j + 1} ^ {y} - S _ {j} ^ {y} S _ {j + 1} ^ {x} \right) \right| \Psi_ {0} \right\rangle = 0.
\tag{5.53}
\end{equation*}

展开 (5.52) 中的余弦，并用不等式

\begin{equation*}
\left\langle S _ {j} ^ {x} S _ {j + 1} ^ {x} + S _ {j} ^ {y} S _ {j + 1} ^ {y} \right\rangle \leq S ^ {2} ,
\tag{5.54}
\end{equation*}

得

\begin{equation*}
\langle \Psi_ {1} | \mathcal {H} | \Psi_ {1} \rangle - E _ {0} \leq \frac {2 \pi^ {2} J S ^ {2}}{\mathcal {N}} + \mathcal {O} ( \mathcal {N} ^ {- 3} ) ,
\tag{5.55}
\end{equation*}

命题 5.9 得证。证毕。

在热力学极限下，基态与这个低激发态的混合可能破缺 $\mathcal{H}$ 的某个对称性。例如 Majumdar--Ghosh 模型\footnote{一个次近邻相互作用特别强的 Heisenberg 模型，见 8.2.1 小节。}的基态与第一激发态都是破缺格子对称性的二聚体组态的叠加。另一种情形发生在自旋 1/2 的近邻 Heisenberg 模型中：des Cloizeaux 与 Pearson 由精确解发现磁振子激发是无能隙的。这些以格子动量 $\mathbf{k}$ 标记的激发，能量为

\begin{equation*}
\omega_ {\mathbf {k}} = ( \pi / 2 ) \left| \sin k \right| ,
\tag{5.56}
\end{equation*}

在热力学极限下于 $k \rightarrow 0$ 与 $k \rightarrow \pi$ 处趋于零。这些激发并非自旋波[见 (11.40)]，因为它们不是绕破缺自旋对称性的基态的小涨落。

与半奇整数自旋链相反，第 15 章我们将学到整数自旋链的激发谱有能隙，称为 Haldane 能隙。历史上，des Cloizeaux 与 Pearson 的激发早于 Haldane 对整数自旋链能隙的预言。而如今，半经典方法使整数情形在概念上反而比半奇整数情形更简单：后者的无能隙被理解为拓扑 Berry 相位之间的量子干涉，我们将在 15.1 节回到这个话题。

\begin{exercises}

\item 考虑无穷程二分哈密顿量

\begin{equation*}
H = + | J | \sum_ {i \in A,\, j \in B} \mathbf {S} _ {i} \cdot \mathbf {S} _ {j} ,
\tag{5.57}
\end{equation*}

求基态与基态能量。验证 Marshall 定理 5.1。

\item 循反铁磁体 Marshall 定理的证明，证明对铁磁体 (5.36)，在 Ising 组态（固定磁化 $M$）中展开得到有限系数 $f_{\alpha}$ 同号。

\item 求长程铁磁体

\begin{equation*}
\tilde {H} = - | J | \sum_ {ij} \mathbf {S} _ {i} \cdot \mathbf {S} _ {j}
\tag{5.58}
\end{equation*}

基态的能量与自旋，并把结果与习题 1 结合证明推论 5.6。

\end{exercises}

\begin{chapbib}
\item Marshall 的原始论文：
  W. Marshall, Proc. R. Soc. London Ser. A \textbf{232}, 48 (1955)。
\item 本章取材于：
  E. Lieb and D. C. Mattis, J. Math. Phys. \textbf{3}, 749 (1962)，
  他们推广并强化了 Marshall 定理。
\item 半奇整数自旋链无能隙性的证明：
  E. Lieb, T. D. Schultz, and D. C. Mattis, Ann. Phys. \textbf{16}, 407 (1961)。
\item 近邻反铁磁体的磁振子激发（由 Bethe 解计算）：
  J. des Cloizeaux and J. J. Pearson, Phys. Rev. \textbf{128}, 2131 (1962)。
\end{chapbib}
